\chapter{RTL을 보고 면적을 예측하는 직관}

\begin{summarybox}[면적을 읽는 핵심]
FPGA 면적은 하나의 gate count가 아니다. LUT, FF, carry, DSP, BRAM, slice와 routing
수요를 따로 본다. 먼저 bit 폭, 병렬 instance 수, pipeline stage, MUX input 수,
memory port를 손으로 세고, 합성 보고서에서 constant propagation과 packing이 만든
차이를 확인한다.
\end{summarybox}

\section{LUT, FF, DSP와 BRAM을 하나의 숫자로 더하지 않는다}

LUT 한 개, FF 한 개, DSP 한 개와 BRAM 한 tile은 서로 교환 가능한 같은 면적
단위가 아니다. 특정 device에서 각각의 가용량과 물리적 위치가 다르다. 따라서
면적 표에는 최소한 다음을 분리해 쓴다.

\begin{itemize}
  \item combinational logic를 나타내는 LUT와 carry primitive
  \item state와 pipeline을 나타내는 FF, control set과 slice packing
  \item multiplication을 담당하는 DSP48E1
  \item storage를 담당하는 RAMB36/RAMB18 tile과 LUTRAM
  \item LUT/FF 수만으로 보이지 않는 slice, congestion과 clock/control 배선
\end{itemize}

``LUT를 줄였다''는 문장과 ``FPGA 전체 면적이 줄었다''는 문장은 같은 뜻이 아니다.
DSP를 추가해 LUT를 줄였을 수도 있고, BRAM 폭 낭비가 늘었을 수도 있다. 목표
device에서 부족한 자원이 무엇인지까지 말해야 설계 판단이 된다.

\section{폭과 병렬도에서 시작하는 손계산}

폭 $W$의 bitwise function을 $P$개 병렬로 두면 첫 근사는 $W\times P$개의
1-bit output function이다. 각 output이 LUT6 한 개에 들어간다면 LUT 상한은
대략 이 값에서 시작한다.

\begin{lstlisting}
generate
  for (g = 0; g < 4; g = g + 1)
    assign y[g*23 +: 23] = a[g*23 +: 23] ^ b[g*23 +: 23];
endgenerate
\end{lstlisting}

이 구조는 23-bit XOR lane 네 개이므로 최대 92개의 1-bit LUT function을
떠올린다. 실제 LUT 수는 다른 논리와의 결합, dual-output packing과 사용되지 않는
bit 제거로 작아질 수 있다.

register는 저장해야 하는 live bit를 센다.

\[
N_{FF,rough}=\sum_{stage} (data\ bits+valid\ bits+mode\ bits+address\ bits).
\]

23-bit data 네 lane을 여섯 stage 동안 모두 유지하면 data만 552 FF다. 그러나
한 stage 이후에 필요 없는 mode별 rail까지 함께 끌고 가면 실제 정보량보다 FF가
늘어난다. Unified NTT에서 radix-3에만 필요한 D lane의 tail valid를 제거한 것은
이 live-range 관점의 사례다.

\section{MUX 면적은 output 폭과 선택 함수로 센다}

$W$-bit 2:1 MUX는 각 output bit마다 세 입력 함수가 필요하므로 isolated mapping의
첫 근사는 $W$ LUT다. $W$-bit 4:1 MUX도 각 bit가 여섯 입력 함수여서 LUT6 한
단계에 들어갈 수 있다. 실제 23-bit 실험에서 둘 다 23 LUT였다.

8:1 이상에서는 각 output bit가 LUT 하나의 입력 범위를 넘는다. LUT tree와
MUXF7/MUXF8을 생각하고, 선택 단계가 늘수록 logic뿐 아니라 select net의 fanout도
늘어난다. 서로 다른 산술 결과를 고르는 MUX라면 MUX 비용 외에 산술기 복제 비용도
센다.

\begin{lstlisting}
// One shared adder: operand MUX + one carry chain
assign shared_y = a + (mode ? b0 : b1);

// Parallel adders: two carry chains + result MUX
assign parallel_y = mode ? (a + b0) : (a + b1);
\end{lstlisting}

첫 구조는 area에 유리할 가능성이 있고, 두 번째는 MUX와 add를 병렬 진행시켜
timing에 유리할 가능성이 있다. 합성기가 공통 연산을 다시 공유할 수도 있으므로
hierarchy와 schematic에서 실제 구조를 확인한다.

\section{Adder, comparator와 correction의 면적}

$W$-bit add/subtract는 대략 $W$개의 LUT-side function과 $\lceil W/4\rceil$개의
CARRY4 구간에서 시작한다. magnitude comparator도 bit 우선순위 또는 borrow를
판단해야 하므로 carry 성격의 구조가 될 수 있다. equality comparator는 bit별
XNOR와 reduction 성격에 가깝다.

다음 modular correction을 연산자 하나로 세면 면적을 과소평가하기 쉽다.

\begin{lstlisting}
assign y = (x >= Q) ? (x - Q) : x;
\end{lstlisting}

구조적으로 comparator, subtractor, result MUX가 필요하다. 상수 $Q$의 bit pattern과
range proof를 이용해 비교와 subtract가 공유되거나 일부 bit가 제거될 수도 있다.
이번 isolated 23-bit 실험에서 일반 subtract는 LUT 23/CARRY4 6, 두 단계
modular subtract는 LUT 46/CARRY4 12였다. 이 차이는 correction이 ``짧은 if문''이
아니라 두 번째 carry propagation을 만들 수 있음을 보여 준다.

\section{Pipeline의 면적은 data만 세면 부족하다}

pipeline 한 stage를 추가할 때에는 data FF 외에 다음 bit들이 함께 이동한다.

\begin{itemize}
  \item result가 유효한지를 나타내는 valid
  \item 연산 종류와 modulus를 나타내는 mode
  \item writeback address, bank, side, port tag
  \item branch별 bypass와 exception 상태
\end{itemize}

4개의 23-bit lane에 1-bit valid 하나만 추가한다고 생각했지만, 실제 interface가
각 lane별 valid와 address를 요구하면 비용이 더 크다. 반대로 phase FSM이 mode를
한 transaction 동안 고정한다면 mode를 모든 stage에 반복 저장하지 않아도 되는지
검토할 수 있다. 이때 기능 계약을 assertion으로 남겨야 한다.

\section{Memory는 bit 수보다 shape와 port가 중요하다}

논리 저장량은 $Depth\times Width$이지만 BRAM tile 수는 primitive의 지원 폭,
depth와 port 구성에 따라 양자화된다. 36 Kb보다 작은 memory도 width나 independent
port 요구 때문에 RAMB36 하나를 사용할 수 있다. 두 read와 한 write를 같은
cycle에 요구하면 단순 dual-port primitive만으로 가능한지, bank 복제나 time
multiplexing이 필요한지 먼저 본다.

\begin{summarybox}[BRAM 면적을 예상할 때 적을 항목]
\begin{enumerate}
  \item logical depth와 word width
  \item synchronous read latency와 read-during-write semantics
  \item read/write port 수와 같은 주소 collision 처리
  \item 병렬 lane을 위한 banking 또는 replication
  \item mode별 memory를 각각 둘지 word packing으로 합칠지
  \item RAMB36과 RAMB18 조합에서 남는 width/depth 공간
\end{enumerate}
\end{summarybox}

Unified NTT의 twiddle ROM은 mode별 ROM을 모두 instantiate하고 MUX로 고르는 대신,
32-bit word 안에 narrow value 두 개를 pack한 2048-word image로 통합했다. Falcon
forward/inverse table까지 포함하면서 twiddle storage는 1.5 tile에서 2 tile로
0.5 늘었지만 coefficient bank는 7.5 tile에서 6 tile로 줄었다. 부분별 증감을
분리해 설명해야 전체 1.5-tile 감소의 원인을 정확히 전달할 수 있다.

\section{자원 공유의 손익 계산}

배타적으로 동작하는 두 operator를 공유하면 큰 arithmetic unit 하나를 줄일 수
있다. 대신 operand MUX, mode decode, result alignment register와 output MUX가
생긴다.

\[
\Delta A_{share}
=-A_{operator}+A_{input\ mux}+A_{control}+A_{alignment}+A_{output\ mux}.
\]

DSP처럼 operator 자체가 큰 hard block이면 공유 이득이 클 수 있다. 9-bit adder처럼
작은 operator를 공유하려다 23-bit MUX와 여러 alignment register가 붙으면 오히려
면적이 늘 수 있다. 공유 여부는 RTL의 연산자 개수가 아니라 post-synthesis
hierarchy와 critical path를 같이 보고 결정한다.

최종 \rtlfile{unifiedMUL}은 ML-DSA와 네 narrow-prime Montgomery mode 사이에서
DSP 두 개, arithmetic unit과 stage register를 공유한다. mode가 transaction 동안
고정되고 매 cycle 새 입력을 받는다는 계약 덕분에, operator를 시간 직렬화해
throughput을 잃지 않고 같은 six-stage pipeline을 사용할 수 있었다.

\section{Packing 때문에 LUT/FF 수와 slice 수가 달라진다}

같은 200 LUT와 200 FF라도 LUT output을 바로 같은 slice의 FF가 받는 구조와,
서로 다른 reset/enable을 가진 FF가 먼 net을 받는 구조는 slice 사용과 timing이
다르다. 그래서 다음 항목을 함께 본다.

\begin{itemize}
  \item slice LUTs와 slice registers의 절대 수
  \item occupied slice와 LUT-FF pair utilization
  \item control set 수와 high-fanout control
  \item LUTRAM/SRL 때문에 필요한 SLICEM 분포
  \item carry chain, DSP, BRAM 주변 placement와 congestion
\end{itemize}

post-synthesis와 post-route LUT 수가 약간 달라지는 것도 이상이 아니다. physical
optimization이 logic replication, LUT combining 또는 trimming을 수행할 수 있다.
반드시 같은 report stage끼리 비교한다.

\section{면적 예측 예제: 23-bit modular subtract 두 lane}

다음 구조가 두 lane 병렬이라고 가정한다.

\begin{lstlisting}
wire [22:0] d0 = a0 - b0;
wire [22:0] y0 = d0[22] ? d0 + Q : d0;
wire [22:0] d1 = a1 - b1;
wire [22:0] y1 = d1[22] ? d1 + Q : d1;
\end{lstlisting}

손계산 순서는 다음과 같다.

\begin{enumerate}
  \item lane당 subtract와 correction add로 carry chain 두 개를 예상한다.
  \item 23-bit 결과 MUX와 sign/borrow 판정을 확인한다.
  \item 두 lane이므로 조합 자원 근사를 두 배 한다.
  \item output register가 있다면 46 FF와 공통 valid/address FF를 더한다.
  \item $a,b$의 범위를 이용해 correction 한 번이면 충분한지 증명한다.
  \item 합성 후 CARRY4와 LUT 수, duplicated logic, trimmed MSB를 확인한다.
\end{enumerate}

이번 operator experiment의 modular subtract 한 개가 LUT 46/CARRY4 12였으므로,
같은 isolated 조건의 두 lane은 LUT 92/CARRY4 24를 첫 기준으로 잡을 수 있다.
실제 NTT hierarchy에서는 upstream/downstream logic과 상수 전파 때문에 달라진다.

\section{면적을 줄였다고 판단하는 절차}

\begin{checklistbox}
\begin{enumerate}
  \item 변경 전후 top, parameter, device, synthesis directive와 report stage를 맞춘다.
  \item 삭제한 live bit, operator, memory word와 port를 RTL에서 직접 센다.
  \item LUT, FF, slice, DSP, BRAM을 각각 비교한다.
  \item hierarchy별로 감소 위치가 예상한 module인지 확인한다.
  \item timing 악화, cycle 증가와 throughput 저하가 없는지 함께 측정한다.
  \item 여러 자원을 하나의 임의 area 숫자로 합치지 않는다.
\end{enumerate}
\end{checklistbox}
